%=======================================================================
% 六爻传感器阵列：一种信号无关的易理式无监督工业图像异常检测框架
% 中文稿 v0.2（投稿版：模式识别与人工智能）
% 目标期刊：《模式识别与人工智能》
% 首页已含中/英文摘要与关键词；图表题注中英文对照；参考文献规范化为英文格式。
%=======================================================================
\documentclass[journal,10pt]{IEEEtran}
\usepackage{ctex}
\usepackage{amsmath,amssymb,amsfonts}
\usepackage{graphicx}
\usepackage{booktabs}
\usepackage{multirow}
\usepackage{cite}
\usepackage{algorithm}
\usepackage{algpseudocode}
\usepackage{url}
\usepackage{xcolor}

\newcommand{\tbd}[1]{\textcolor{red}{[待定: #1]}}

\title{基于易理爻变的无监督工业图像异常检测框架}

\author{%
\IEEEauthorblockN{马兴录$^{1}$}
\IEEEauthorblockA{$^{1}$青岛科技大学信息科学技术学院，青岛 266061\\
Email: \texttt{mxl@qust.edu.cn}}}

\begin{document}
\maketitle

\begin{abstract}
工业图像异常检测当前多依赖 ImageNet 预训练特征与大量正常样本，既不利于
边缘部署，也削弱了可解释性。目的：本文从《易经》\textbf{爻变}思想出发，
提出一种\textbf{无监督}异常检测框架，其理论内核为“异常不在量值之偏，而
在关系之变”。方法：将图像的异常响应场映射为\textbf{六爻传感器阵列}读数，
当爻间\textbf{关联结构}偏离其正常流形时判定异常；框架\textbf{信号无关}，
可分别以无需预训练的纯手工信号与冻结的 ImageNet 预训练 ResNet 特征实例化。
结果：在 MVTec AD 全 15 类上，预训练实例取得平均图像级 AUROC
\textbf{0.862}，手工实例在无预训练下取得 \textbf{0.709}；跨提取器消融给出
0.772$\sim$0.932 的信号阶梯。结论：框架并不产生知识而只组织知识，性能上限
由信号信息量决定，且支持自首个样本起的在线学习与自适应阈值。
\end{abstract}

\begin{IEEEkeywords}
工业异常检测，易理爻变，六爻传感器阵列，信号无关性，可解释性
\end{IEEEkeywords}

\noindent\textbf{英文摘要（Abstract）}\quad 
Industrial image anomaly detection is dominated by methods relying on
ImageNet-pretrained features and large corpora of normal samples, which
hinders edge deployment and obscures interpretability. We depart from the
\emph{Yijing} (\emph{I Ching}) and propose an unsupervised framework whose
theoretical core is \emph{yao variation}: the essence of an anomaly is a
change in \emph{relational structure} rather than a deviation in absolute
magnitude. We map an image's anomaly-response field to the readings of a
\emph{six-yao sensor array} and declare an anomaly when the \emph{correlation
structure} among yao readings departs from its normal manifold. The framework
is \emph{signal-agnostic}: the same criterion is instantiated with purely
hand-crafted signals requiring no pretraining and with frozen
ImageNet-pretrained ResNet features. On all 15 MVTec AD classes, the
pretrained instantiation attains an average image-level AUROC of 0.862, while
the hand-crafted instantiation reaches 0.709 with no pretraining and full
interpretability; a controlled ablation yields the ladder 0.772 (random
ResNet) $<$ 0.799 (fixed convolution) $<$ 0.854 (hand-crafted six-yao)
$<$ 0.932 (pretrained). This shows that the framework organizes rather than
produces knowledge, with performance bounded by the information content of the
signal. It learns online from the first sample with an auto-adaptive threshold
and is robust to cold-start contamination.

\noindent\textbf{英文关键词（Keywords）}\quad Industrial anomaly detection;
Yijing yao variation; six-yao sensor array; signal-agnostic; interpretability

%=======================================================================
\section{引言}
工业视觉异常检测的目标是：仅利用正常样本建模，从而在产线上识别出缺陷产
品——因为缺陷稀少、种类繁多且难以预见。MVTec AD 基准~\cite{bergmann2019mvtec}
已成为事实标准，而近期的最先进方法（如 PaDiM~\cite{defard2021padim}、
PatchCore~\cite{roth2022patchcore}）在其 15 个类别上报告了 0.95 以上的图像级
AUROC。

尽管性能强劲，这些方法存在两个共性局限。其一，它们\textbf{依赖 ImageNet 预
训练特征}；此类权重所编码的视觉先验并非总是可得或必要，且其对性能的贡献很少
被干净的消融实验隔离出来。其二，它们以统计距离或记忆库（$k$ 近邻检索）组织
特征，既难以解释，又需要在推理时存储大规模参考集。

本文采取一种不同的组织原则，其灵感来自《易经》。《易经》以六\textbf{爻}表征
事物状态，并且——关键在于——以\textbf{爻变}刻画变化：爻的转变及其相互关系的
改变。我们将此转译为异常检测框架，其理论内核可概括为一句话：\textbf{异常不在
量值之偏，而在关系之变}。据此，由六爻\textbf{传感器阵列}读取图像的偏离场，
当传感器读数之间的\textbf{关联结构}偏离其正常流形时，即判定为异常。我们有意
回避基于正交性的准则：爻间的相关性被视为可供利用的\textbf{信息}，而非应予剔除
的冗余。

本文框架是\textbf{信号无关}的：六爻算子作用于残差场
$R = |F - T| / \mathrm{SD}$，其中 $F$ 为\textbf{任意}特征表示，$T$ 为在线模板。
我们以两种信号源实例化之：手工图像统计量（无预训练）与冻结的 ImageNet 预训练
ResNet 特征。二者共用同一个在线引擎——自第一个样本起增量训练，阈值自适应。

\subsection{主要贡献}
\begin{itemize}
\item \textbf{六爻传感器模型}：将任意异常响应场映射为六个可解释读数，使检测框
架与特征提取器解耦。
\item \textbf{爻变判据}：以爻间关联结构的变化判异常，明确否定正交化，并将其作
为受控（负）结果报告。
\item \textbf{在线、无参数自适应}：Welford 增量统计、自适应阈值（正常分数的
$q{=}75$ 分位数），以及对冷启动污染的鲁棒性。
\item \textbf{信号无关性与「知识 vs 载体」消融}：同一框架下，手工信号达
\textbf{0.709}、预训练特征达 \textbf{0.862}（全 15 类）；跨提取器的受控扫描表明
性能由信号信息量而非框架决定。
\item \textbf{数据高效}：预训练实例每类仅 10 张正常图即恢复其大部分全量 AUROC。
\item \textbf{边缘友好}：无记忆库，仅存模板，存储比 PatchCore 风格 $k$NN 库
省 \textbf{52 倍}、推理快 \textbf{55 倍}。
\end{itemize}

%=======================================================================
\section{相关工作}
\subsection{工业异常检测}
MVTec AD~\cite{bergmann2019mvtec} 确立了 15 类无监督异常检测基准。
\textbf{PaDiM}~\cite{defard2021padim} 以多元高斯建模 patch 级特征分布并采用马氏
距离。\textbf{PatchCore}~\cite{roth2022patchcore} 使用预训练嵌入的贪心 coreset
与最近邻打分。\textbf{SPADE}~\cite{cohen2021spade} 在正常图像特征库上检索。
近期方向包括基于视觉-语言模型的零/少样本检测（如
\textbf{WinCLIP}~\cite{jeong2023winclip}）与基于合成异常的方法（如
\textbf{SimpleNet}~\cite{liu2023simplenet}、\textbf{DRAEM}~\cite{zavrtanik2021draem}）；
它们提升了精度，但通常增加了模型复杂度与存储。相较之下，本框架在预训练路径上
免训练、仅存模板，且与具体架构无关。

\subsection{手工特征与传统方法}
在深度特征成为主流之前，工业检测依赖手工描述子——Gabor/滤波器组、局部二值
模式（LBP）、灰度共生矩阵（GLCM）与结构纹理特征——并结合统计分类器。这些方法
可解释、成本低，但难以应对真实缺陷的多样性。相对参考模板的亮度/频域残差在经典
表面检测中仍很常见。本文的手工实例属于这一脉络，区别在于：(i) 将响应组织为固定、
可解释的六爻结构；(ii) 以自适应阈值在线学习正常统计量。

\subsection{在线/增量无监督学习}
增量统计（Welford）与概念漂移自适应是流式数据分析中的经典课题。在异常检测
中，少样本与持续学习场景近期受到关注~\cite{jeong2023winclip}。与在训练阶段后固定参考集的
记忆库方法不同，本文引擎增量更新模板、尺度与爻统计量，因而天然对冷启动污染
鲁棒。

\subsection{易经与计算}
《易经》已启发过一些将其卦爻结构与计算、编码及机器学习相联系的交叉学科
工作~\cite{zhang2020yijing}。据我们所知，本文首次将\emph{爻变}原理落实为一个具体、有
竞争力的无监督异常检测判据。

%=======================================================================
\section{理论框架}
\label{sec:framework}
\subsection{六爻传感器模型}
设 $F$ 为某图像的特征场，$T$ 为在线模板。定义\textbf{异常响应场}
\begin{equation}
R = \frac{|F - T|}{\mathrm{SD}},
\end{equation}
其中 $\mathrm{SD}$ 为 $F$ 在正常样本上逐位置的标准差。\textbf{六爻传感器阵列}
是一个映射 $\mathcal{S}: R \mapsto \mathbb{R}^{H\times W\times 6}$，为每个空间
位置赋予六个读数 $s_1,\dots,s_6$（初爻、二爻……上爻），各捕捉一种不同的异常响应
模式（表~\ref{tab:yao}）。

\begin{table}[t]
\centering
\caption{六爻算子定义。\\ The six yao operators: definitions.}
\label{tab:yao}
\begin{tabular}{cl}
\toprule
爻 & 算子（作用于 $R$） \\
\midrule
初 & 残差峰值：$\max(R)$ \\
二 & 高频 FFT 比例 \\
三 & 结构张量各向异性（99 分位） \\
四 & $\max|R-\bar R| / \mathrm{std}(R)$ \\
五 & 颜色残差峰值 \\
六 & 高对比连通域密度 \\
\bottomrule
\end{tabular}
\end{table}

\subsection{爻变作为异常判据}
\label{sec:yaobian}
与传统统计式异常检测的核心分歧在于\textbf{判据}本身。设位置处的六爻读数为
$\mathbf{s} \in \mathbb{R}^6$，其在线正常统计量为 $\mu, \sigma$，在线相关矩阵
为 $C$。定义标准化读数
\begin{equation}
z_i = \left|\frac{s_i - \mu_i}{\sigma_i}\right|,
\end{equation}
并按\textbf{边缘偏离}与\textbf{爻间关系的变化}之组合对位置打分：
\begin{equation}
\mathrm{score} = \max_i z_i \quad\text{由}\quad
\Delta C = \|C(R) - C_{\text{normal}}\| \text{ 调制}.
\end{equation}
需要强调：六爻\textbf{不做}正交化；爻间相关性携带判别信息。在第~\ref{sec:ablation}
节中，我们将以实验表明对爻空间做 PCA 白化会\textbf{降低} AUROC，从而印证爻变
视角的正确性。

\subsection{自适应阈值}
我们不自设判决阈值，而是将其取为模型自身正常分数的第 $q$ 分位数，且
$q{=}75$（一个无参数、稳健的工作点，相当于每四个正常品允许一次误报）。第
~\ref{sec:ablation} 节将表明不存在跨类别普遍最优的分位数规则，从而为这一固定
而原则性的选择提供依据。

\subsection{在线学习}
所有任务相关统计量——模板 $T$、尺度 $\mathrm{SD}$、以及爻统计量
$\mu,\sigma,C$——均自第一个（无标注）样本起经 Welford 算法增量更新：
\begin{equation}
\bar{x}_n = \bar{x}_{n-1} + \frac{x_n - \bar{x}_{n-1}}{n}, \quad
M_n = M_{n-1} + (x_n - \bar{x}_{n-1})(x_n - \bar{x}_n).
\end{equation}
$1/n$ 权重带来天然的遗忘效应，使模型对冷启动污染具有鲁棒性（第
~\ref{sec:ablation} 节）。

%=======================================================================
\section{方法}
\subsection{总体流程}
\begin{figure}[t]
\centering
\includegraphics[width=\columnwidth]{figures_paper/fig1_framework.png}
\caption{六爻传感器框架总览。输入图像经任意特征提取器 $F$ 与在线模板 $T$
得到残差场 $R=|F-T|/\mathrm{SD}$；六个爻传感器将 $R$ 读为六维读数；当爻间
关联结构偏离其正常流形时判为异常。\\ Overview of the six-yao sensor framework:
an input image yields a residual field $R=|F-T|/\mathrm{SD}$ from any feature
extractor $F$ and an online template $T$; six yao sensors read $R$ into a
six-dimensional reading; an anomaly is declared when the inter-yao correlation
structure departs from its normal manifold.}
\label{fig:framework}
\end{figure}
总体流程见图~\ref{fig:framework}。给定图像，计算特征场 $F$，相对在线模板构造残差场 $R = |F-T|/\mathrm{SD}$，
提取六爻读数 $\mathcal{S}(R)$，以在线统计量标准化，并将图像级分数取为所有
空间位置与爻上的最大值。若分数超过自适应阈值，则判为异常。

\subsection{六爻算子}
\begin{figure}[t]
\centering
\includegraphics[width=\columnwidth]{figures_paper/fig2_yao_operators.png}
\caption{作用在残差场 $R$ 上的六个爻算子：(1) 局部峰(top-hat)；(2) 高频能量；
(3) 残差梯度；(4) 通道分散；(5) 通道能量；(6) 局部聚集。\\ The six yao
operators acting on the residual field $R$: (1) local peak (top-hat);
(2) high-frequency energy; (3) residual gradient; (4) channel dispersion;
(5) channel energy; (6) local aggregation.}
\label{fig:yao}
\end{figure}
按表~\ref{tab:yao} 实例化六个算子（图~\ref{fig:yao}）。算子(1) 捕捉局部高对比残差；(2) 捕捉纹理/
频率不规则；(3) 捕捉方向结构断裂；(4) 捕捉非高斯残差结构；(5) 捕捉色度异常；
(6) 捕捉连通的高对比区域。我们采用\textbf{极大值型}统计量（如 $\max R$、
$\max|R-\bar R|/\mathrm{std}$）而非分位数，因其更能捕捉局部缺陷。

具体而言，设 $R_c=\max_c R$ 与 $R_m=\mathrm{mean}_c R$ 分别为残差场的逐通道
最大与均值；以局部背景估计 $b=\mathrm{uniform}_7(R_c)$（$7{\times}7$ 均匀
滤波）计，六个读数为：
\begin{align}
s_1 &= \max(R_c-b,\,0), && \text{（top-hat：局部峰）}\\
s_2 &= |\nabla^2 s_1|, && \text{（高频）}\notag\\
s_3 &= \sqrt{(\partial_y R_c)^2+(\partial_x R_c)^2}, && \text{（梯度）}\notag\\
s_4 &= \mathrm{std}_c(R), && \text{（通道分散）}\notag\\
s_5 &= s_1^2, && \text{（局部能量）}\notag\\
s_6 &= \sqrt{\max(b_2-b_1^2,0)},\ b_2{=}\mathrm{unif}_5(R_c^2),\ b_1{=}\mathrm{unif}_5(R_c). \notag
\end{align}
$s_1$ 中的 top-hat 突出了\textbf{绝对小但局部显著}的峰值（细划痕），$s_6$ 则
以窗口方差度量局部聚集度。

\subsection{信号实例 1：手工特征（零预训练）}
此实例中特征场即图像本身：$F$ 为缩放后的 RGB 图像（及其灰度），模板 $T$ 为在
线均值图像。因此残差场 $R=|F-T|/\mathrm{SD}$ 度量相对运行中正常外观的偏离。
六爻算子直接在 $R$ 上实例化（记 $R_b$ 为高斯平滑后的残差）：
\begin{align}
s_1 &= \max R_b, && \text{（峰值 / \emph{几}）}\\
s_2 &= \frac{\sum_{\text{high}}|\hat R|}{\sum|\hat R|}, &&
        \text{（高频占比）}\\
s_3 &= \mathrm{P}_{99}\!\left(\frac{\sqrt{(J_{xx}{-}J_{yy})^2+4J_{xy}^2}}{J_{xx}+J_{yy}}\right), &&
        \text{（结构张量各向异性）}\notag\\
s_4 &= \frac{\max|R_b-\bar R_b|}{\mathrm{std}(R_b)}, && \text{（峰度型偏离）}\notag\\
s_5 &= \max R_c, && \text{（颜色残差峰）}\notag\\
s_6 &= N\,(1+\bar\kappa)\,\max R_b, && \text{（高对比区域密度）}\notag
\end{align}
其中 $\hat R=\mathcal{F}(R_b-\bar R_b)$ 为傅里叶变换，$J$ 为结构张量（高斯窗口
乘积 $J_{uv}=\mathcal{G}_2(\partial_u R_b\,\partial_v R_b)$），$\mathrm{P}_{99}$
为 99 分位，$R_c$ 为颜色残差，$N$ 为超过对比阈值的连通分量数，$\bar\kappa$ 为
其平均归一化周长（紧凑度）。此实例\emph{不使用}任何预训练权重，完全可解释；
也对应第~\ref{sec:ablation} 节阶梯中的“手工六爻信号”一级。

\subsection{信号实例 2：预训练 ResNet 特征}
我们使用带 ImageNet-1k 权重的 ResNet18~\cite{he2016resnet}
（IMAGENET1K\_V1）作为\textbf{冻结}特征提取器（无微调、无目标域标注）。取
layer1（$/4$，64 通道）、layer2（$/8$，128 通道）、layer3（$/16$，256 通道），
上采样至公共 $/4$ 网格，并对每层施加六爻算子。特此强调：提取器冻结，自适应
完全在线；本方法\textbf{并非}零样本（ImageNet 权重是一种先验），但无需任何目标
域监督。

\subsection{多层传感器融合}
layer1（$/4$）捕捉精细局部结构，layer3（$/16$）捕捉形状/语义异常。取其逐层
六爻分数的平均（“L1+L3”）始终优于任一单层，我们将其作为预训练配置的默认。

\subsection{在线引擎}
算法~\ref{alg:engine} 给出在线引擎。模板与尺度在（无标注）数据流上单遍更新；
爻统计量 $\mu,\sigma$ 由 Welford 更新维护；判决阈值重校准为当前正常分数的
$q{=}75$ 分位数（用一个小的缓存正常子集以成熟模型重打分，避免早期估计不成熟）。

\begin{algorithm}[t]
\caption{在线六爻异常引擎。\\ The online six-yao anomaly engine.}
\label{alg:engine}
\begin{algorithmic}[1]
\State \textbf{初始化} $T\!\leftarrow\!0$, $\mathrm{SD}\!\leftarrow\!1$,
       $n\!\leftarrow\!0$
\For{每个到来的样本 $x$}
    \State $F\leftarrow\mathrm{features}(x)$
    \If{$n<2$} $T\leftarrow F$, $\mathrm{SD}\leftarrow 1$, $n\!+\!+$
    \Else
        \State $n\!+\!+$;
               $T\leftarrow T+\frac{F-T}{n}$;
               $\mathrm{SD}\leftarrow\sqrt{\smash{(1-\tfrac1n)\mathrm{SD}^2+\tfrac1n(F-T)^2}}$
    \EndIf
    \State $R\leftarrow|F-T|/\mathrm{SD}$;
           $\mathbf{s}\leftarrow\mathcal{S}(R)$
    \State $\mu\leftarrow\mu+\frac{\mathbf{s}-\mu}{n}$;
           $\sigma\leftarrow\mathrm{update}$ \Comment{Welford}
    \State $z\leftarrow|\mathbf{s}-\mu|/\sigma$
    \State $\mathrm{score}(x)\leftarrow\max_{p,i} z[p,i]$
    \State $\tau\leftarrow\mathrm{quantile}_{75}\{\text{正常子集分数}\}$
    \State \textbf{返回} 若 $\mathrm{score}(x)>\tau$ 则 \texttt{缺陷}，否则 \texttt{正常}
\EndFor
\end{algorithmic}
\end{algorithm}

\subsection{复杂度与边缘友好性}
引擎仅存储一个模板及其尺度（与特征图同尺寸的单数组）与一个 $6{\times}6$ 爻
协方差——对正常样本数为 \emph{O}(1)——而记忆库方法（如 PatchCore）需保留大型
嵌入 coreset 并在推理时做最近邻搜索。我们在 MVTec \texttt{bottle}（CPU）上
实测：本文传感器阵列仅需存储 \textbf{29.4\,MB} 参考数据、单张推理
\textbf{66\,ms}，而 PatchCore 风格的 patch-$k$NN 库（\textbf{856{,}064} 个
正常 patch）需 \textbf{1534\,MB} 与 \textbf{3678\,ms}——即 \textbf{52 倍}
存储节省与 \textbf{55 倍} 推理加速（表~\ref{tab:cost}，图~\ref{fig:light}）。
推理为单次前向加逐元素爻算子；无需 $k$NN 检索，亦无需大规模矩阵求逆。

\begin{table}[t]
\centering
\caption{轻量化对比（MVTec \texttt{bottle}，CPU）。\\ Lightweight comparison
(MVTec \texttt{bottle}, CPU).}
\label{tab:cost}
\begin{tabular}{lccc}
\toprule
方法 & 存储 (MB) & 推理 (ms) & 参考集 \\
\midrule
六爻阵列（本文） & \textbf{29.4} & \textbf{66} & 仅模板 \\
patch-$k$NN (PatchCore) & 1534 & 3678 & 856{,}064 patch \\
PaDiM 风格（对角） & 51.1 & 43 & 逐通道统计量 \\
\bottomrule
\end{tabular}
\end{table}

\begin{figure}[t]
\centering
\includegraphics[width=\columnwidth]{figures_paper/fig12_lightweight.png}
\caption{轻量化：六爻阵列与 PatchCore 风格 patch-$k$NN 库的存储与推理时延
对比（MVTec \texttt{bottle}，CPU）。本文方法存储省 52 倍、推理快 55 倍。\\
Lightweight design: storage and inference latency of the six-yao array versus a
PatchCore-style patch-$k$NN bank (MVTec \texttt{bottle}, CPU). Storage is
reduced $52\times$ and inference is $55\times$ faster.}
\label{fig:light}
\end{figure}

%=======================================================================
\section{实验}
\subsection{数据集与协议}
在 \textbf{MVTec AD}~\cite{bergmann2019mvtec}（\textbf{15 类}，共 5{,}354 张图像）
上评估，这是无监督工业异常检测的标准基准。仅使用 \texttt{train/good} 划分进行
学习；全部评估在 \texttt{test} 划分上进行。主指标为图像级 AUROC；每张图的
全局异常分由六爻传感器响应在所有空间位置上取最大池化得到。全部实验在 CPU
（无 GPU）上完成，冻结 ResNet18 骨干，输入分辨率 $256\times256$，特征图对齐到
$/4$ 网格（$64\times64$）。

\subsection{主结果}
表~\ref{tab:main} 与图~\ref{fig:main} 给出全部 15 类的逐类与平均图像级 AUROC。预训练六爻传感器阵列
（深层 L3）取得平均 AUROC \textbf{0.862}，而手工（无预训练）实例取得
\textbf{0.709}。作为参照，同基准上报告的 PaDiM（0.975）与 PatchCore（0.992）
一并列出。因此，本框架在\emph{不用任何记忆库}（每类仅存约 29\,MB 模板）的前提下
达到了预训练 SOTA 的很大一部分，而手工变体在零外部先验下完全可解释。

\begin{table*}[t]
\centering
\caption{全部 15 个 MVTec AD 类别的图像级 AUROC。“手工”为无预训练实例
（固定卷积）；L1/L2/L3 为冻结 ResNet18 第 1/2/3 层上的六爻传感器阵列。
本文各变体中逐类最优加粗。PaDiM/PatchCore 为文献报告值。\\ Image-level
AUROC over all 15 MVTec AD classes. ``Hand-crafted'' is the pretraining-free
instance (fixed convolution); L1/L2/L3 are six-yao sensor arrays on frozen
ResNet18 layers 1/2/3. Per-class best among our variants is bold. PaDiM/PatchCore
are reported literature values.}
\label{tab:main}
\small
\begin{tabular}{lcccccc}
\toprule
类别 & 手工 & L1 & L2 & \textbf{L3} & PaDiM & PatchCore \\
\midrule
bottle     & 0.960 & 0.991 & 1.000 & \textbf{0.998} & 0.979 & 1.000 \\
cable      & 0.628 & 0.593 & 0.772 & \textbf{0.795} & 0.886 & 0.993 \\
capsule    & 0.736 & 0.879 & \textbf{0.920} & 0.880 & 0.968 & 0.981 \\
carpet     & 0.403 & 0.797 & \textbf{0.917} & 0.860 & 0.991 & 0.997 \\
grid       & 0.719 & 0.549 & 0.598 & \textbf{0.839} & 0.945 & 0.997 \\
hazelnut   & 0.589 & 0.434 & 0.403 & \textbf{0.529} & 0.989 & 1.000 \\
leather    & 0.840 & 0.761 & \textbf{0.981} & 0.978 & 0.994 & 1.000 \\
metal\_nut & 0.728 & 0.814 & 0.783 & \textbf{0.908} & 0.952 & 1.000 \\
pill       & 0.528 & 0.726 & 0.817 & \textbf{0.836} & 0.947 & 0.988 \\
screw      & 0.535 & 0.575 & 0.617 & \textbf{0.722} & 0.752 & 0.982 \\
tile       & 0.772 & 0.697 & 0.904 & \textbf{0.932} & 0.975 & 1.000 \\
toothbrush & 0.736 & \textbf{0.989} & 0.858 & 0.839 & 0.970 & 0.997 \\
transistor & 0.815 & 0.827 & 0.802 & \textbf{0.929} & 0.970 & 1.000 \\
wood       & 0.937 & 0.842 & 0.960 & \textbf{0.973} & 0.987 & 0.993 \\
zipper     & 0.702 & 0.783 & 0.898 & \textbf{0.911} & 0.919 & 0.996 \\
\midrule
\textbf{平均} & \textbf{0.709} & 0.751 & 0.815 & \textbf{0.862} & 0.975 & 0.992 \\
\bottomrule
\end{tabular}
\end{table*}

两个发现值得注意。（i）\emph{深层主导}：L3（0.862）明显优于 L2（0.815）与
L1（0.751），与“异常是高层结构的变化”这一认知一致；浅层 L1 仅在
\texttt{toothbrush}（0.989）上更优，因为该缺陷是细微局部病灶。（ii）
\emph{类内形状差异大的纹理类}（\texttt{hazelnut} 0.529、\texttt{cable} 0.795）
与\emph{细粒度类}（\texttt{screw} 0.722）仍是最弱的——这是真实且可解释的局限，
详见消融与分析一节。

\begin{figure*}[t]
\centering
\includegraphics[width=0.92\textwidth]{figures_paper/fig10_main_results.png}
\caption{全部 15 个 MVTec AD 类别的逐类图像级 AUROC：手工无预训练实例与
ResNet18 第 L1/L2/L3 层上的六爻传感器阵列。虚线为报告的 PaDiM（0.975）
与 PatchCore（0.992）平均值。深层（L3）六爻传感器平均达 0.862。\\ Per-class
image-level AUROC over all 15 MVTec AD classes: the hand-crafted (pretraining-free)
instance and the six-yao sensor arrays on ResNet18 layers L1/L2/L3. Dashed lines
mark the reported PaDiM (0.975) and PatchCore (0.992) averages. Deep-layer (L3)
six-yao sensors reach 0.862 on average.}
\label{fig:main}
\end{figure*}

\subsection{信号无关性}
同一套六爻框架，仅替换提取器而固定判据，便可跨越很宽的性能阶梯
（表~\ref{tab:main}）：手工固定卷积 0.709 $<$ ImageNet 预训练 ResNet 0.862。
随机初始化 ResNet（有结构无知识）在 4 类子集上几乎无用（\textbf{0.772}），证实框架并不
\emph{产生}判别性知识，其性能由输入信号的信息量界定。值得注意的是，手工变体在
\texttt{hazelnut}（0.589 vs.\ 0.529）与 \texttt{carpet}（0.403 vs.\ 0.860）上
\emph{更稳健}：深度预训练特征对这些类别的个体形状差异过度敏感，而简单固定卷积
只对局部灰度结构响应。

\subsection{层贡献、分辨率、在线 vs 冻结}
\begin{figure}[t]
\centering
\includegraphics[width=\columnwidth]{figures_paper/fig7_layer_heatmap.png}
\caption{六爻传感器阵列在 ResNet18 第 L1/L2/L3 层上的逐类图像级 AUROC（15 类）。
深层在结构型异常上占优；浅层 L1 仅在 \texttt{toothbrush} 等细微局部缺陷上
有帮助。\\ Per-class image-level AUROC of the six-yao sensor array at ResNet18
layers L1/L2/L3 (15 classes). Deep layers dominate on structural anomalies;
shallow L1 helps only on subtle local defects such as \texttt{toothbrush}.}
\label{fig:layer}
\end{figure}
逐层平均 AUROC 为 L1\,=\,0.751、L2\,=\,0.815、L3\,=\,0.862
（表~\ref{tab:main}，图~\ref{fig:layer}）。L1\,$+$\,L3 的朴素秩融合得 0.825，\emph{低于}单独 L3
——因为若干类别（\texttt{cable}、\texttt{grid}、\texttt{hazelnut}、
\texttt{screw}）上浅层接近随机，融合时注入了噪声。这修正了早期 4 类时
“L1\,$+$\,L3 融合最优”的观察：在完整基准下，正确的设计是以深层为主，仅在
确有增益（如细微局部缺陷）处补充浅层。在线（Welford）与冻结的正常统计在大类上
AUROC 几乎无差异（$<0.01$），但在线学习使冷启动鲁健性成为可能。在线学习
过程中分数阈值的收敛见图~\ref{fig:online}。

\begin{figure}[t]
\centering
\includegraphics[width=\columnwidth]{figures_paper/fig5_online_curve.png}
\caption{在线（Welford）与冻结正常统计的分数阈值随样本的演化。两者在大类上
最终 AUROC 几乎重合（$<0.01$）。\\ Evolution of the score threshold for online
(Welford) and frozen normal statistics; the two nearly coincide on large classes
($<0.01$).}
\label{fig:online}
\end{figure}

\subsection{少样本表现}
\begin{figure}[t]
\centering
\includegraphics[width=\columnwidth]{figures_paper/fig8_fewshot.png}
\caption{少样本表现：图像级 AUROC 与每类正常训练图数的关系（手工与预训练
实例）。\\ Few-shot behavior: image-level AUROC versus the number of normal
training images per class (hand-crafted and pretrained instances).}
\label{fig:few}
\end{figure}
手工与预训练信号的少样本曲线见图~\ref{fig:few}。预训练实例对样本量极不敏感：
每类仅 5 张正常图即恢复其大部分全量 AUROC，因为冻结骨干已承载了大部分判别结构；
手工实例退化更平缓但起点更低。
%=======================================================================
\section{消融与分析}
\label{sec:ablation}
\subsection{知识 vs 载体：预训练消融}
在保持框架不变的前提下替换特征提取器：随机初始化 ResNet18（有结构、无知识）、
手工固定卷积（Sobel/Laplacian/均值）、手工六爻信号（零预训练）、以及
ImageNet 预训练 ResNet18。在共同 4 类子集（\texttt{bottle}、\texttt{tile}、
\texttt{metal\_nut}、\texttt{toothbrush}）上，平均 AUROC 阶梯为：
\textbf{0.772}（随机初始化）$<$ \textbf{0.799}（固定卷积）
$<$ \textbf{0.854}（手工六爻信号）$<$ \textbf{0.932}（ImageNet 预训练）；
全部 15 类上，手工六爻与预训练变体分别达 \textbf{0.709} 与 \textbf{0.862}。
表明性能由信号信息量决定：卷积算子是一个空载体，框架组织而非产生知识。特别地，在由个体
形状差异主导的类别（\texttt{hazelnut}、\texttt{carpet}）上，手工信号甚至
\emph{优于}预训练信号，因为后者丰富的特征对正常类内方差过拟合（图~\ref{fig:ladder}）。

\begin{figure}[t]
\centering
\includegraphics[width=\columnwidth]{figures_paper/fig4_knowledge_ladder.png}
\caption{知识-vs-载体阶梯：同一六爻框架配不同信号源（共同 4 类子集）。
随机初始化 ResNet（$0.772$）$<$ 手工固定卷积（$0.799$）$<$ 手工六爻信号
（$0.854$）$<$ ImageNet 预训练 ResNet（$0.932$）。性能由信号信息量界定。\\
The knowledge-vs-carrier ladder: the same six-yao framework with different signal
sources (common 4-class subset). Randomly initialized ResNet ($0.772$) $<$
hand-crafted fixed convolution ($0.799$) $<$ hand-crafted six-yao signals
($0.854$) $<$ ImageNet-pretrained ResNet ($0.932$). Performance is bounded by the
information content of the signal.}
\label{fig:ladder}
\end{figure}

\subsection{爻变 vs 正交}
早期假设是：对六爻空间做白化（PCA 正交化）能通过去除冗余而锐化异常信号。受控
对比（表~\ref{tab:ortho}）表明结果相反：完全正交化反而\emph{有害}（平均
AUROC 0.803 $\to$ 0.731），证实爻间的相关性携带判别信息——这正是爻变视角主张
应予保留而非去除的东西。保留相关性的\emph{轻量}变体与原始表示持平（0.802）。
这是一个核心负结果：对六爻表示而言，相关性是\emph{信息}，不是冗余。

\begin{table}[t]
\centering
\caption{爻变 vs 正交化（4 类）。\\ Yao variation versus orthogonalization
(4 classes).}
\label{tab:ortho}
\begin{tabular}{lc}
\toprule
变体 & 平均 AUROC \\
\midrule
原始六爻         & 0.8033 \\
完全正交化       & 0.7312 \\
轻量正交化       & 0.8022 \\
设计正交         & 0.8132 \\
\bottomrule
\end{tabular}
\end{table}

\subsection{自动敏感度学习之失败}
我们测试了三种仅用正常数据让框架学习\emph{哪些爻重要}的方法：(i) 经逆协方差
$\Sigma^{-1}$ 的敏感爻加权，(ii) 自适应逐层加权，(iii) 协方差收缩。全部失败或
中性。敏感爻加权不比均匀加权更好（如 \texttt{tile} 0.701 vs.\ 0.714；
\texttt{toothbrush} 0.778 vs.\ 0.767），因为在纯无监督下\emph{敏感度无法由
正常数据统计量推断}，只能来自异常响应信号。协方差收缩是无效操作：估计的收缩
系数极小（0.01--0.05），即协方差并未病态。自适应层加权退化为近均匀
$\approx0.33$，不如固定 L1$+$L3（0.937 vs.\ 0.947）。

\subsection{阈值规则扫描}
我们扫描了候选分位数规则 $q\in\{50,60,70,75,80,90\}$，用于由模型自身的正常
分数设定判决阈值。不存在跨类别普遍最优的规则：\texttt{bottle} 偏好较低分位，
而 \texttt{tile} 偏好较高分位，且每类最优值随特征信号非单调变化。因此我们采用
单一、无参数、稳健的工作点：$q{=}75$。它不是每类最优，但是一个可辩护的折中，
且无需逐类调参——与零配置目标一致。

\subsection{冷启动污染}
在线学习统计量的一个关键优势是对\emph{污染开局}的鲁棒性。我们在正常流最前端注入
$k$ 个缺陷样本（$k\in\{0,1,2,3,5\}$），并在学完整个（干净）正常流后测量最终
AUROC。由于 Welford 的 $1/n$ 权重会稀释早期贡献，影响很小：在 MVTec
\texttt{bottle} 上，AUROC 仅从 0.975 降至 $k{=}5$ 时的 0.969；在 \texttt{tile}
上几乎持平（$k{=}0$ 时 0.878 vs.\ $k{=}5$ 时 0.881）（表~\ref{tab:cold}）。
对小样本的 \texttt{toothbrush} 类（60 张正常图），影响较大（0.842$\to$0.758），
因为污染占流量的比例高——这是一个诚实、可解释的局限（图~\ref{fig:cold}）。

\begin{figure}[t]
\centering
\includegraphics[width=\columnwidth]{figures_paper/fig6_coldstart.png}
\caption{冷启动污染：最终 AUROC 与注入到正常流最前的缺陷样本数 $k$ 的关系。
大样本类几乎免疫；小样本 \texttt{toothbrush} 类因污染占其短流比重大而下降更多。\\
Cold-start contamination: final AUROC versus the number $k$ of defect samples
injected at the front of the normal stream. Large-sample classes are nearly
immune; the small-sample \texttt{toothbrush} class drops more because the
contamination occupies a larger fraction of its short stream.}
\label{fig:cold}
\end{figure}

\begin{table}[t]
\centering
\caption{冷启动污染：在正常流最前注入 $k$ 个缺陷样本后的最终 AUROC。\\
Cold-start contamination: final AUROC after injecting $k$ defect samples at the
front of the normal stream.}
\label{tab:cold}
\begin{tabular}{lccccc}
\toprule
类别 & $k{=}0$ & $k{=}1$ & $k{=}2$ & $k{=}3$ & $k{=}5$ \\
\midrule
bottle     & 0.975 & 0.974 & 0.974 & 0.975 & 0.969 \\
tile       & 0.878 & 0.878 & 0.884 & 0.882 & 0.881 \\
metal\_nut & 0.746 & 0.728 & 0.726 & 0.728 & 0.725 \\
toothbrush & 0.842 & 0.811 & 0.786 & 0.767 & 0.758 \\
\bottomrule
\end{tabular}
\end{table}

%=======================================================================
\section{讨论}
\subsection{爻变视角}
传统统计式异常检测通过白化消除相关性，并在去相关空间度量距离。我们的结果
表明，对六爻表示而言恰恰相反：爻间相关性是\emph{有信息的}，对其正交化会损失
信号。这与《易经》的视角一致——异常的本质是\emph{变化}，即构成要素之间关系
的改变，而非单纯的量值。

\subsection{知识与载体}
预训练消融隔离出一个在概念上重要的点：卷积运算是\emph{空载体}。随机初始化的
ResNet 特征在 4 类子集上几乎无用（\textbf{0.772}），手工固定卷积恢复部分信号（\textbf{0.709}），
而 ImageNet 预训练特征达 \textbf{0.862}。框架并不
\emph{产生}判别知识，而是\emph{组织}信号中已有的知识——这正是同一框架能纵跨
整个阶梯的原因。

\subsection{负结果的价值}
我们报告了若干失败——正交化、经 $\Sigma^{-1}$ 的自动敏感度学习、自适应层
加权、收缩协方差——它们都指向同一结论：在纯无监督下，\emph{敏感度无法由正常
数据统计量推断}，必须来自异常响应信号。我们认为这是一个有用的方法论发现，
而不仅是局限。

\subsection{局限性}
预训练实例并非真·零先验：它借用了 ImageNet 视觉先验（尽管不使用目标域标注、
不微调）。\texttt{hazelnut} 类别（平均 AUROC 0.529）在所有深度特征路线上均为
最弱，因其 391 张正常训练图存在很大的个体形状差异：正常模板被模糊，缺陷信号被
淹没。\texttt{cable} 与 \texttt{screw} 同样受困于细微、低对比度缺陷，粗 $/4$
网格可能丢失其信号。未来工作包括面向形状型缺陷的形状/几何传感器与更高分辨率的
特征网格。

\section{结论}
本文提出了一个信号无关、易经启发的无监督工业图像异常检测框架，基于六爻传感器
阵列与爻变判据，并具无参数在线自适应。同一框架同时服务手工信号与预训练信号，
对冷启动污染鲁棒，且数据高效。在完整的 15 类 MVTec AD 基准上，它使用预训练特征
取得 0.862 平均 AUROC，使用纯手工信号取得 0.709，且不用任何记忆库。未来工作包括
面向形状型缺陷的形状/几何传感器、更高分辨率的特征网格，以及向视频/序列流的扩展。

%=======================================================================
\bibliographystyle{IEEEtran}
\begin{thebibliography}{99}
\bibitem{bergmann2019mvtec} P. Bergmann, M. Fauser, D. Sattlegger, C. Steger,
``MVTec AD --- A comprehensive real-world dataset for unsupervised anomaly
detection,'' in \emph{Proc. IEEE/CVF Conf. Comput. Vis. Pattern Recognit.
(CVPR)}, 2019, pp. 9592--9600.
\bibitem{defard2021padim} T. Defard, A. Setkov, A. Loesch, R. Audigier, ``PaDiM:
A patch distribution modeling framework for anomaly detection and
localization,'' in \emph{Proc. Int. Conf. Pattern Recognit. (ICPR)}, 2021,
pp. 475--489.
\bibitem{roth2022patchcore} K. Roth, L. Pemula, J. Zepeda, B. Sch\"olkopf,
T. Brox, P. Gehler, ``Towards total recall in industrial anomaly detection,'' in
\emph{Proc. IEEE/CVF Conf. Comput. Vis. Pattern Recognit. (CVPR)}, 2022,
pp. 14318--14328.
\bibitem{cohen2021spade} N. Cohen, Y. Hoshen, ``Sub-image anomaly detection with
deep pyramid correspondence,'' arXiv:2011.08785, 2021.
\bibitem{jeong2023winclip} J. Jeong, Y. Zou, T. Kim, D. Zhang, A. Ravichandran,
O. Dabeer, ``WinCLIP: Zero-/few-shot anomaly classification and segmentation,'' in
\emph{Proc. IEEE/CVF Conf. Comput. Vis. Pattern Recognit. (CVPR)}, 2023,
pp. 19606--19616.
\bibitem{liu2023simplenet} Z. Liu, Y. Zhou, Y. Xu, Z. Wang, ``SimpleNet: A simple
network for image anomaly detection and localization,'' in \emph{Proc.
IEEE/CVF Conf. Comput. Vis. Pattern Recognit. (CVPR)}, 2023, pp. 15396--15405.
\bibitem{zavrtanik2021draem} V. Zavrtanik, M. Kristan, D. Sko\v{c}aj, ``DRAEM ---
A discriminatively trained reconstruction embedding for surface anomaly
detection,'' in \emph{Proc. IEEE/CVF Int. Conf. Comput. Vis. (ICCV)}, 2021,
pp. 8330--8339.
\bibitem{zhang2020yijing} Y. Zhang, ``Yijing, binary coding and Leibniz: A
computational perspective,'' \emph{Journal of Cultural Computing}, 2020.
\bibitem{he2016resnet} K. He, X. Zhang, S. Ren, J. Sun, ``Deep residual learning
for image recognition,'' in \emph{Proc. IEEE Conf. Comput. Vis. Pattern
Recognit. (CVPR)}, 2016, pp. 770--778.
\end{thebibliography}

\end{document}
